\section{Model Usage}
% \begin{itemize}
%     \item How could your model be used?
%     \item Why could your model be used?
%     \item By whom could your model be used?
%     \item What needs to happen so users can make a request and get a usable response?
%     \item What happens before your model runs?
%     \item What happens after your model makes a prediction?
%     \item How will you turn your model's prediction into a usable answer for the user?
%     \item \textbf{Required:} At least one user group and one use case.
%     \item \textbf{Required:} Something that happens before and something that happens after your model.
% \end{itemize}

As hypothesised in the introduction and based on preliminary domain knowledge, the model could be used for theft and suspicious activity detection, purchase confirmation based on recent motion behaviour, and other lightweight security checks requiring identity confirmation. While these are all practical applications, they are feasible and realistic due to the prospective system being relatively small, capable of running on a device and having access to live motion data. The input of acceleration and rotation data and the output of an onset/offset verification have a wide range of applications, and the model itself will likely be part of a larger onboard security system. The great thing is that anyone with a smartphone, especially those who want an additional passive layer of security, can use the model. Since our initial goal is to create a working service with just an API endpoint, this can be used for less live and onboard security, similar to software such as FindMy and Life360, which perform passive and possibly interval checks. We don't want to livestream biometric data, so ideally we want edge deployment.

In order for the model to function, there would first be a short enrolment period during which the user would walk for a few minutes in order to create an identification embedding or other enrolment, such as a centroid of features. After that, live data would be streamed in windows based on step timing, as described in the original paper. This data would be interpolated and resampled to handle jitter, and then compared against the known owner embedding using cosine similarity or another comparison method. Most of the time, the system would remain idle after the model. If it detects a potential threat, the app could trigger an alert, message or notification on another device. The final output is essentially a yes/no decision, so the user-facing result would probably be a notification on another device. In order to support this, we would require SMS, email, bot-based messaging or a server that can receive the signal and send a ping to another device. This could act as theft detection for children with phones, with parents receiving notifications when their child gives the device to an untrusted adult or it is stolen.